\begin{center}
{\LARGE\bfseries \ReportTitle\par}
\vspace{0.8em}
{\large \StudentAuthors\par}
\vspace{0.4em}
{\small \SchoolAffiliation\par}
\vspace{0.4em}
{\small \SupervisingTeachers\par}
\end{center}

\section*{Abstract}
\noindent A detector can recognize that a monitored time series changed without
knowing whether the change is local to one target or shared across a comparison
network. That distinction determines which follow-up evidence is relevant, yet
real monitoring records rarely provide reliable scope labels. We formulate
\emph{selective cross-site auditing} as an information-constrained problem: at
a stated scope-error tolerance, which cases can an observation channel answer,
and when should it abstain?

\noindent MetaShift-Bench holds the target path exactly fixed across paired
local and shared synthetic worlds and varies only donor-side participation. Its
finite predeclared policy set gives target-only information no positive
held-out coverage below 50\% tolerated error. Comparative information reaches
\VFiveComparativeFrontierTwentyPercent{} coverage only at a 20\% answered-case
error tolerance ($\alpha=.20$), with \VFiveComparativeRiskTwentyPercent{}
observed conditional error. The structural interval-separation certificate
abstains on \VFiveCertificateAbstentionPercent{} of matched pairs and is neither
a population-risk guarantee nor a deployable real-network certificate.

\noindent We prove the target-only impossibility statement under matched,
balanced target laws, derive the availability-normalized partial-scope residual
gap, and give a sufficient condition for structural abstention. Across
\TotalAnchors{} reported AQS metadata anchors, \CompleteComparisons{} permitted
a complete comparison; \SupportedCandidates{} were supported audit candidates,
\NotSupportedEvents{} were not supported, and \InconclusiveEvents{} were
inconclusive. These are audit dispositions, not physical faults, causal
mechanisms, or ground-truth labels. In an independent stable-regime benchmark,
cross-site methods outscore target-only baselines, but MetaShift has no
confidence-supported aggregate advantage over standard synthetic control. The
central contribution is therefore not estimator superiority, but an evaluation
framework in which an audit answers a scope question only when its comparison
evidence can support that answer.

\noindent\textbf{Keywords:} local-versus-shared change; time-series auditing;
selective prediction; cross-site comparison; structural abstention.
